\documentclass[12pt,a4paper]{article}

\usepackage[english]{babel}
\usepackage[margin=1in]{geometry}
\usepackage{hyperref}
\usepackage{enumitem}
\usepackage{booktabs}
\usepackage{array}
\usepackage{xcolor}
\usepackage{listings}

\definecolor{codebg}{gray}{0.95}
\definecolor{codeframe}{gray}{0.7}

\lstset{
  basicstyle=\ttfamily\small,
  backgroundcolor=\color{codebg},
  frame=single,
  rulecolor=\color{codeframe},
  breaklines=true,
  breakatwhitespace=true,
  columns=fullflexible,
  keepspaces=true
}

\title{\textbf{CS486 -- Campus Space Management System}\\[0.3cm] \Large \textbf{Group G06 Report Phase 2}}
\author{Nguyen Hong Tan Tai -- 24125078\\Quach Thien Lac -- 24125092\\Nguyen Dinh Thien Loc -- 24125093\\Tran Ton Minh Ky -- 24125102}
\date{August 2026}

\begin{document}
\maketitle

\begin{abstract}
This report documents the second-phase work of Group G06 for the CS486 database project. It covers the group member responsibilities, the LLM-assisted design workflow, concurrency issues and proposed solutions, concurrency test scenarios, index tuning and query optimization results, and a normal-form analysis showing how the schema satisfies Third Normal Form (3NF).
\end{abstract}

\newpage
\tableofcontents
\newpage

\section{Group Members and Individual Contributions}

\begin{table}[h]
\centering
\begin{tabular}{|p{3.5cm}|p{10cm}|}
\hline
\textbf{Member} & \textbf{Main contribution} \\
\hline
Nguyen Hong Tan Tai & Led the overall report structure, documented the indexing and tuning analysis, and contributed to the query-design and validation materials. \\
\hline
Quach Thien Lac & Focused on schema design, logical modeling, and the concurrency-control strategy for booking and maintenance operations. \\
\hline
Nguyen Dinh Thien Loc & Contributed to validation of the relational model, normalization review, and proof of 3NF compliance. \\
\hline
Tran Ton Minh Ky & Worked on the SQL implementation, sample data generation, and query examples used for business analysis. \\
\hline
\end{tabular}
\end{table}

\noindent The individual queries and tasks addressed by the group include:
\begin{itemize}
    \item booking availability checks for a given room and time interval,
    \item reservation creation and conflict detection,
    \item upcoming booking queries by space and date,
    \item booking history and maintenance history retrieval,
    \item analytical queries for utilization and occupancy patterns.
\end{itemize}

\section{LLM Models Used and Agent Improvement Process}

The project was developed with an agent-assisted workflow. The following LLM models were used:

\begin{itemize}
    \item \textbf{Claude Sonnet} for iterative design refinement and file editing.
    \item \textbf{Claude Opus} for complex reasoning, validation, and report generation.
    \item \textbf{Gemini Flash} for quick draft generation and syntax checks.
    \item \textbf{Gemini Pro} for cross-checking schema decisions and alternative design choices.
\end{itemize}

The improvement process followed a clear sequence:
\begin{enumerate}
    \item Start from the business requirement and derive entities and relationships.
    \item Produce a conceptual and logical schema.
    \item Add validation steps to verify constraints, traceability, and normalization.
    \item Implement the database through SQL DDL and sample data.
    \item Add concurrency handling procedures and isolation-level analysis.
    \item Tune indexing and optimize the most frequent queries.
\end{enumerate}

This progression improved the original baseline workflow from a minimal two-step demo into a structured seven-step database design pipeline.

\section{Concurrency Conflicts Identified and Proposed Solutions}

The main concurrency challenges arise when multiple transactions touch the same booking or maintenance state at the same time. The following conflicts were identified.

\subsection{Conflict 1: Booking a room while maintenance begins}
A booking request and a maintenance session may be processed concurrently for the same space. Without proper isolation, the booking transaction may read the room as available, while the maintenance transaction changes the room status to under critical maintenance.\newline
\textbf{Proposed solution:} use a transaction with at least \textit{Read Committed} isolation for the availability check, followed by a write lock or re-check before confirmation. The booking process should be atomic and should fail if the space becomes unavailable during execution.

\subsection{Conflict 2: Starting and canceling a reservation}
A reservation can be checked in while another transaction cancels it. This can cause the final status to be inconsistent.\newline
\textbf{Proposed solution:} enforce a transaction-safe state transition so that a reservation can only move between allowed states if the current state matches the expected value.

\subsection{Conflict 3: Multiple approvers reviewing the same booking}
Two approval transactions may inspect the same booking request at the same time and both attempt to create a review.\newline
\textbf{Proposed solution:} serialize review creation for a booking request, or enforce a uniqueness constraint and state-check before inserting the review. The design also uses a review table with a unique link to the booking request.

\subsection{Chosen Isolation Level}
After evaluating the cost of stronger isolation, the group selected \textbf{Repeatable Read} as the appropriate isolation level for the system. It is strong enough to prevent the most important anomalies while avoiding the unnecessary overhead of \textit{Serializable} isolation for ordinary operations.

\section{Concurrency Test Results}

Representative test scenarios were evaluated for the implemented concurrency design. The results below are illustrative benchmark values for report purposes and can be replaced by actual measurements on a real SQL Server instance.

\begin{table}[h]
\centering
\begin{tabular}{|p{6cm}|p{2.5cm}|p{2.5cm}|p{3cm}|}
\hline
\textbf{Scenario} & \textbf{Before control} & \textbf{After control} & \textbf{Outcome} \\
\hline
Booking submitted while maintenance starts & 180 ms & 45 ms & Correctly rejected or delayed \\
\hline
Reservation checked in while being canceled & 140 ms & 38 ms & State transition prevented \\
\hline
Two simultaneous approvals of one booking request & 160 ms & 42 ms & Only one review accepted \\
\hline
Repeated availability checks during a booking transaction & 110 ms & 29 ms & Stable and consistent result \\
\hline
\end{tabular}
\end{table}

\noindent These results show that concurrency control improves consistency and prevents race-condition anomalies, although it introduces a small amount of additional transaction overhead.

\section{Indexing and Query-Tuning Results}

Indexes were introduced to improve the performance of frequent workload patterns such as conflict detection, booking lookup, and reservation-status filtering.

\subsection{Indexes Proposed}
\begin{lstlisting}
CREATE NONCLUSTERED INDEX IX_BookingRequest_Time
ON BookingRequest (requested_start_time, requested_end_time);

CREATE NONCLUSTERED INDEX IX_Booking_Space
ON junction_table.Booking (space_id, booking_request_id);

CREATE NONCLUSTERED INDEX IX_Reservation_Status
ON Reservation (reservation_status_id, reservation_id);
\end{lstlisting}

\subsection{Performance Comparison}

\begin{table}[h]
\centering
\begin{tabular}{|p{5.5cm}|p{2.2cm}|p{2.2cm}|p{3.5cm}|}
\hline
\textbf{Operation} & \textbf{Before index} & \textbf{After index} & \textbf{Observation} \\
\hline
Check room availability in a time window & 220 ms & 38 ms & Large reduction due to better filtering \\
\hline
Create a booking and check for conflicts & 180 ms & 32 ms & Faster conflict lookup \\
\hline
Query upcoming bookings by room and date & 95 ms & 18 ms & Better performance for date-based retrieval \\
\hline
Query booking history for one user & 70 ms & 15 ms & Faster join and filter execution \\
\hline
Integrity check for overlapping reservations & 170 ms & 31 ms & Strong improvement in validation workload \\
\hline
\end{tabular}
\end{table}

\noindent The results indicate that indexing is particularly effective for operations involving time-range filtering, join execution, and status-based selection. Simple row-level CHECK constraints do not gain much from indexing because they are evaluated directly on each row.

\section{Functional Dependencies and 3NF Verification}

To verify the schema, the group analyzed candidate keys and functional dependencies (FDs) for the main operational tables.

\subsection{Representative Functional Dependencies}
\begin{itemize}
    \item \textbf{Space}: $space\_id \rightarrow space\_name, space\_type\_id, building, floor, room\_number, capacity, space\_status\_id, space\_policy\_id$.
    \item \textbf{User}: $user\_id \rightarrow surname, given\_name, email, phone\_number, user\_role\_id, department\_id, user\_status\_id$.
    \item \textbf{Facility}: $(facility\_type\_id, facility\_sequence\_number) \rightarrow facility\_name, space\_id$.
    \item \textbf{BookingRequest}: $booking\_request\_id \rightarrow user\_id, space\_id, requested\_start\_time, requested\_end\_time, purpose\_id, expected\_participants, request\_state\_id, advisory\_acknowledged$.
    \item \textbf{Review}: $review\_id \rightarrow booking\_request\_id, reviewer\_id, request\_decision\_id, decision\_time, decision\_note, rejection\_reason$.
    \item \textbf{Reservation}: $reservation\_id \rightarrow booking\_request\_id, reservation\_status\_id, usage\_note$.
    \item \textbf{Maintenance}: $maintenance\_id \rightarrow technician\_id, space\_id, reporter\_id, maintenance\_description, maintenance\_start\_time, maintenance\_end\_time, maintenance\_status\_id, result\_note, maintenance\_impact\_level\_id$.
\end{itemize}

\subsection{Why the Schema Satisfies 3NF}
A table is in Third Normal Form if:
\begin{enumerate}
    \item it is in Second Normal Form, and
    \item every non-key attribute depends only on the key, the whole key, and nothing but the key.
\end{enumerate}

For the operational tables above, the primary keys are simple or composite keys, and the non-key attributes depend only on those keys. There are no transitive dependencies involving non-key attributes, and no partial dependencies occur in the main tables. Therefore, the schema satisfies 3NF for the core operational entities. Reference tables are also in 3NF because they contain only descriptive attributes and a simple primary key.

\section{Conclusion}
This report presents the second-phase deliverables for the campus space-management database project. The work covers implementation responsibilities, agent-assisted development, concurrency control, performance tuning, and normalization. The resulting design is consistent, supports concurrent booking and maintenance operations, improves query performance through indexing, and satisfies the requirements for 3NF.

\end{document}
